\expandafter\ifx\csname CAPETINCLUDEMODE\endcsname\relax
\documentclass[11pt,a4paper]{article}
% preamble.tex — packages et commandes communes a toutes les fiches CAPET SII IE
\input{glyphtounicode}
\pdfgentounicode=1
\usepackage{cmap}
\usepackage[T1]{fontenc}
\usepackage[utf8]{inputenc}
\usepackage[french]{babel}
% babel[french] active a \begin{document} des raccourcis de ponctuation (:;!?)
% qui cassent la syntaxe to[...] de circuitikz (erreur "+ or - expected") ;
% on les desactive, ce qui ne change que l'espacement fin avant ces signes.
\AtBeginDocument{\shorthandoff{:;!?}}
\usepackage[a4paper, top=1.6cm, bottom=1.6cm, left=1.4cm, right=1.4cm]{geometry}
\setlength{\headheight}{14.5pt}
\usepackage{amsmath, amssymb, amsfonts}
\usepackage{siunitx}
\sisetup{locale=FR, output-decimal-marker={,}, inter-unit-product=\ensuremath{{}\cdot{}},
         group-minimum-digits=4}
% Unites non fournies par siunitx, declarees ici une fois pour toutes (fiches + exercices).
% Les gardes \ifdefined evitent tout conflit avec une redeclaration locale restee dans un
% fichier : ne pas retirer les gardes, ne pas redeclarer ces unites dans les fichiers.
\DeclareSIUnit{\jour}{j}
\ifdefined\var\else\DeclareSIUnit{\var}{var}\fi   % puissance reactive
\ifdefined\vad\else\DeclareSIUnit{\vad}{VAD}\fi   % puissance deformante
\ifdefined\tr\else\DeclareSIUnit{\tr}{tr}\fi      % tour (tr/min)
\usepackage{array, booktabs, multirow, tabularx}
\usepackage{xcolor}
\usepackage[most, breakable]{tcolorbox}
\usepackage{multicol}
\usepackage{enumitem}
\usepackage{mathtools}
\usepackage{lmodern}
% Refonte visuelle 2026 (direction "moderne coloré") : Lato pour le texte,
% maths sans-serif assorties via newtxsf (charge apres amsmath/amssymb).
\usepackage[default]{lato}
\usepackage{newtxsf}
\usepackage{microtype}
\usepackage{fancyhdr}
\usepackage{tikz}
\usetikzlibrary{positioning,arrows.meta}
\usepackage{circuitikz}
\usepackage{lastpage}

% ── Couleurs par domaine (palette modernisée 2026) ────────────────────
\definecolor{colElec}{RGB}{29, 78, 216}        % Électrotechnique — bleu vif
\definecolor{colElecLight}{RGB}{219, 234, 254} % fond bleu clair
\definecolor{colAuto}{RGB}{21, 128, 61}        % Automatique — vert
\definecolor{colMeca}{RGB}{194, 65, 12}        % Mécanique — orange
\definecolor{colTherm}{RGB}{185, 28, 28}       % Thermique — rouge
\definecolor{colQT}{RGB}{232, 244, 233}        % (hérité, plus utilisé par questiontype)
\definecolor{colJury}{RGB}{255, 248, 225}      % fond jury (crème)
\definecolor{colJuryFrame}{RGB}{202, 138, 4}   % bordure/bandeau jury (ambre)
\definecolor{colGrayLight}{RGB}{245, 245, 247} % fond gris neutre
\definecolor{colSeparator}{RGB}{180, 180, 180} % gris séparateur
\definecolor{colExo}{RGB}{8, 145, 178}          % Banque d'exercices — cyan (distinct des fiches)
% ── Encre et gris de la refonte ───────────────────────────────────────
\definecolor{colInk}{RGB}{23, 32, 42}          % texte courant (quasi-noir)
\definecolor{colSlate}{RGB}{83, 98, 118}       % gris ardoise (méta, en-têtes)
\definecolor{colRule}{RGB}{203, 213, 225}      % filets clairs
\definecolor{colJuryBg}{RGB}{255, 248, 225}    % fond crème jury
\definecolor{colPresence}{RGB}{83, 98, 118}    % bloc 5 (présence dans les sujets)
\AtBeginDocument{\color{colInk}}

% ── Style encadré exercice + corrigé (banque d'exercices) ─────────────
\tcbset{
  exoBox/.style={
    breakable,
    enhanced,
    colback=colExo!8!white,
    colframe=colExo,
    colbacktitle=colExo,
    coltitle=white,
    fonttitle=\bfseries\small,
    left=4pt, right=4pt, top=3pt, bottom=3pt,
    before skip=5pt, after skip=5pt,
  }
}

% ── Bascule enonce / corrige (banque d'exercices) ─────────────────────
% \CAPETCORRIGE est defini en ligne de commande (pdflatex "\def\CAPETCORRIGE{1}\input{...}")
% pour compiler la variante corriges d'un meme fichier source ; par defaut (non defini),
% la compilation produit la variante enonces. \CAPETSHOW{enonce}{corrige} choisit la branche.
%
% TROIS MODES (refonte 2026-09-20, manuel de revision) :
%   - aucun drapeau   -> imprime l'enonce            (variante enonces, INCHANGE)
%   - \CAPETCORRIGE   -> imprime le corrige          (variante corriges, INCHANGE)
%   - \CAPETMANUEL    -> imprime l'enonce ET met le corrige de cote pour l'annexe
%
% Le mode manuel stocke chaque corrige dans une macro (\csgdef), sans ecriture
% fichier ni re-tokenisation : le contenu (circuitikz, accents, %) traverse intact.
% L'annexe les rejoue ensuite dans l'ordre via \CAPETANNEXE.
\usepackage{etoolbox}

\newcounter{CAPETexo}          % numero continu d'exercice dans le manuel
\setcounter{CAPETexo}{0}
\newcounter{CAPETannexe}       % compteur de relecture en annexe

% Titre courant de la notion en cours (pose par \CAPETNOTION), sert d'intertitre annexe
\newcommand{\CAPETnotioncode}{}
\newcommand{\CAPETnotiontitre}{}

% \CAPETNOTION{code}{titre} : ouvre une notion dans le manuel (memorise le libelle,
% et marque le rang du premier exercice de la notion pour l'intertitre d'annexe).
\newcommand{\CAPETNOTION}[2]{%
  \renewcommand{\CAPETnotioncode}{#1}%
  \renewcommand{\CAPETnotiontitre}{#2}%
  \csgdef{CAPETnotionstart\the\numexpr\value{CAPETexo}+1\relax}{#1 -- #2}%
}

% ── \degree robuste ───────────────────────────────────────────────────
% Certains fichiers d'exercices utilisent \degree. Ni le preambule ni siunitx
% ne le definissent, et \providecommand local ne suffit pas : sous hyperref,
% \degree existe deja comme commande factice (\ERROR), donc \providecommand
% s'abstient et l'appel echoue a la compilation du manuel (erreur invisible
% en compilation isolee). On le (re)definit donc inconditionnellement ici,
% apres le chargement des packages, pour tout le corpus.
\AtBeginDocument{\providecommand{\degree}{}\renewcommand{\degree}{\ensuremath{^\circ}}}

% Libelle decoratif : les fichiers d'exercices utilisent aussi \CAPETSHOW dans
% leur \fancyhead et leur bloc de titre, pour afficher "(enonces)" ou "(corriges)".
% Ces appels-la ne sont PAS des exercices : en mode manuel ils ne doivent ni
% consommer un numero, ni alimenter l'annexe. On les reconnait a leur premier
% argument, strictement egal au libelle ci-dessous.
\newcommand{\CAPETLIBELLEENONCES}{(\'enonc\'es)}

\ifdefined\CAPETMANUEL
  % --- Mode manuel : enonce en place, corrige differe vers l'annexe ---
  \newcommand{\CAPETSHOW}[2]{%
    \def\CAPETargun{#1}%
    \ifx\CAPETargun\CAPETLIBELLEENONCES
      #1%            appel decoratif : rendu tel quel, hors numerotation
    \else
      \CAPETSHOWEXO{#1}{#2}%
    \fi
  }
  \newcommand{\CAPETSHOWEXO}[2]{%
    \stepcounter{CAPETexo}%
    \csgdef{CAPETcorrige\theCAPETexo}{#2}%
    \edef\CAPETtmpnum{\theCAPETexo}%
    #1%
    \par\vspace{2pt}%
    {\footnotesize\itshape\color{colExo}%
     \hyperref[capet:cor:\CAPETtmpnum]{Corrige de l'exercice \CAPETtmpnum\ $\rightarrow$ p.~\pageref{capet:cor:\CAPETtmpnum}}%
     \label{capet:eno:\CAPETtmpnum}}%
  }
  % \CAPETANNEXE : rejoue tous les corriges captures, dans l'ordre, avec intertitres
  \newcommand{\CAPETANNEXE}{%
    \setcounter{CAPETannexe}{0}%
    \loop\ifnum\value{CAPETannexe}<\value{CAPETexo}%
      \stepcounter{CAPETannexe}%
      \edef\CAPETcurnum{\theCAPETannexe}%
      \ifcsname CAPETnotionstart\CAPETcurnum\endcsname
        \par\medskip
        {\normalsize\bfseries\color{colExo}\csuse{CAPETnotionstart\CAPETcurnum}}%
        \par\nobreak\smallskip
      \fi
      \begin{tcolorbox}[exoBox, breakable,
        title={Corrige de l'exercice \CAPETcurnum\
               \hfill\normalfont\footnotesize
               \hyperref[capet:eno:\CAPETcurnum]{enonce p.~\pageref{capet:eno:\CAPETcurnum}}}]%
      \label{capet:cor:\CAPETcurnum}%
      \csuse{CAPETcorrige\CAPETcurnum}%
      \end{tcolorbox}%
    \repeat
  }
\else
  % --- Modes historiques, comportement strictement inchange ---
  \newcommand{\CAPETSHOW}[2]{\ifdefined\CAPETCORRIGE #2\else #1\fi}
  \newcommand{\CAPETANNEXE}{}
\fi

% ── Difficulté 1 à 3 (étoiles pleines, amssymb) ───────────────────────
\newcommand{\difficulte}[1]{%
  \ifcase#1\relax
  \or $\bigstar$\;niveau 1/3%
  \or $\bigstar\bigstar$\;niveau 2/3%
  \or $\bigstar\bigstar\bigstar$\;niveau 3/3%
  \fi
}

% ── Style encadré extrait de document technique (banque d'exercices) ──
\tcbset{
  docTech/.style={
    breakable,
    enhanced,
    colback=colGrayLight,
    colframe=colSeparator,
    fonttitle=\bfseries\small\itshape,
    coltitle=black,
    colbacktitle=colGrayLight,
    left=4pt, right=4pt, top=3pt, bottom=3pt,
    before skip=4pt, after skip=4pt,
  }
}

% ── Style tcolorbox de section principale ────────────────────────────
% Refonte 2026 : bandeau titre coloré arrondi + fond très pâle, sans cadre.
\tcbset{
  ficheSection/.style={
    breakable, enhanced jigsaw,
    frame hidden, boxrule=0pt, arc=1.8mm,
    colback=#1!4!white,
    colbacktitle=#1,
    coltitle=white,
    fonttitle=\sffamily\bfseries,
    left=7pt, right=7pt, top=5pt, bottom=5pt,
    toptitle=3pt, bottomtitle=3pt,
    before skip=10pt, after skip=8pt,
  }
}

% ── Style encadré question type ──────────────────────────────────────
% Carte blanche, bandeau sombre, liseré très clair.
\tcbset{
  questiontype/.style={
    breakable, enhanced jigsaw,
    colback=white,
    frame hidden, boxrule=0pt, arc=1.8mm,
    colbacktitle=colElec!78!black,
    coltitle=white,
    fonttitle=\sffamily\bfseries\small,
    left=7pt, right=7pt, top=4pt, bottom=4pt,
    toptitle=2.5pt, bottomtitle=2.5pt,
    before skip=8pt, after skip=8pt,
    borderline={0.6pt}{0pt}{colElec!12!white},
  }
}

% ── Style encadré jury ───────────────────────────────────────────────
% Bandeau ambre, fond crème, sans cadre.
\tcbset{
  jury/.style={
    breakable, enhanced jigsaw,
    colback=colJuryBg,
    frame hidden, boxrule=0pt, arc=1.8mm,
    colbacktitle=colJuryFrame,
    coltitle=white,
    fonttitle=\sffamily\bfseries,
    left=7pt, right=7pt, top=4pt, bottom=4pt,
    toptitle=3pt, bottomtitle=3pt,
    before skip=10pt, after skip=8pt,
  }
}

% ── En-tête / pied de page communs (refonte 2026) ─────────────────────
% #1 code fiche, #2 titre court, #3 couleur domaine
\newcommand{\ficheHeaders}[3]{%
  \fancyhead[L]{\small\sffamily{\bfseries\textcolor{#3}{#1}}\hspace{0.7em}\textcolor{colSlate}{#2}}%
  \fancyhead[R]{\small\sffamily\textcolor{colSlate}{CAPET SII IE\, --- \,Session 2026}}%
  \fancyfoot[C]{\small\sffamily\textcolor{colSlate}{\thepage\,/\,\pageref{LastPage}}}%
  \fancyhead[C]{}\fancyfoot[L]{}\fancyfoot[R]{}%
}

% ── Bloc en-tête de fiche (BLOC 0, refonte 2026) ──────────────────────
% #1 code  #2 titre  #3 couleur domaine  #4 domaine  #5 ligne présence  #6 sous-thèmes
\newcommand{\ficheEntete}[6]{%
  \begin{tcolorbox}[enhanced, breakable,
    frame hidden, boxrule=0pt, arc=2.5mm,
    interior style={left color=#3, right color=#3!75!black},
    left=9pt, right=9pt, top=7pt, bottom=7pt,
    before skip=0pt, after skip=6pt]
    {\color{white}\sffamily\bfseries\LARGE #1\quad #2}\\[4pt]
    {\color{#3!22!white}\sffamily\small \MakeUppercase{#4}\quad\textbar\quad #5}\\[1pt]
    {\color{#3!35!white}\sffamily\footnotesize Sous-thèmes : #6}
  \end{tcolorbox}%
}

% ── Mise en page ─────────────────────────────────────────────────────
\pagestyle{fancy}
\fancyhf{}
\renewcommand{\headrulewidth}{0.5pt}
\renewcommand{\footrulewidth}{0pt}
\AtBeginDocument{%
  \renewcommand{\headrule}{{\color{colRule}\hrule width\headwidth height 0.5pt}}%
}

\setlist[itemize]{leftmargin=1.2em, itemsep=1pt, topsep=2pt, parsep=0pt}
\setlist[enumerate]{leftmargin=1.2em, itemsep=3pt, topsep=2pt, parsep=0pt}

\setlength{\parskip}{2pt}
\setlength{\columnsep}{10pt}

\def\CAPETBEGINFILE{\begin{document}}
\def\CAPETENDFILE{\end{document}}
\else
\def\CAPETBEGINFILE{}
\def\CAPETENDFILE{}
\fi

\ficheHeaders{T07}{Hacheurs}{colElec}

\CAPETBEGINFILE

% -----------------------------------------------------------------------
%  BLOC EN-TETE
% -----------------------------------------------------------------------

\ficheEntete{T07}{Hacheurs (1, 2 et 4 quadrants)}{colElec}%
  {électrotechnique - électronique de puissance}%
  {\textbf{Présence :} \textbf{12/13} sujets (2013-2026) - thème \emph{très fréquent}}%
  {hacheur Buck (dévolteur) - hacheur Boost (élévateur) - ondulation courant - hacheur 4 quadrants (pont en H) - entrelacement de bras}

% -----------------------------------------------------------------------
%  1. RAPPELS DE COURS
% -----------------------------------------------------------------------

\begin{tcolorbox}[
  ficheSection=colElec,
  title={1. Rappels de cours},
]

% -- 1.1 Hacheur Buck (devolteur) -------------------------------------
\textbf{1.1 Hacheur Buck - dévolteur (série)}

\smallskip
Principe : l'interrupteur commandé est fermé pendant $t_{on} = \alpha T$ puis ouvert pendant $t_{off} = (1-\alpha)T$.
En conduction continue, la loi de la valeur moyenne donne :
\[
  \langle u_{sortie} \rangle = \alpha \cdot E \qquad \text{avec}\quad \alpha \in [0\,;\,1]
\]

\smallskip
\begin{center}
\begin{tikzpicture}[scale=0.85, font=\footnotesize]
  \draw[->] (0,0) -- (6.3,0) node[right] {$t$};
  \draw[->] (0,-0.3) -- (0,2.2) node[above] {$u_{sortie}(t)$};
  \draw[thick, colElec] (0,1.6) -- (2.4,1.6) -- (2.4,0) -- (6,0);
  \draw[dashed] (0,0.75) -- (6,0.75) node[right] {\scriptsize $\langle u_{sortie}\rangle = \alpha E$};
  \draw[dotted] (2.4,0) -- (2.4,1.6);
  \node[below] at (1.2,0) {\scriptsize $\alpha T$ (ON)};
  \node[below] at (4.2,0) {\scriptsize $(1-\alpha)T$ (OFF)};
  \node[left] at (0,1.6) {\scriptsize $E$};
  \draw[dashed] (6,0) -- (6,-0.3);
  \node[below] at (6,-0.3) {$T$};
\end{tikzpicture}
\end{center}

Ondulation de courant dans l'inductance (pire cas à $\alpha = 0{,}5$) :
\[
  \Delta i_L = \frac{E \cdot \alpha \cdot (1-\alpha)}{L \cdot f}
  \qquad\Rightarrow\qquad
  L_{min} = \frac{E \cdot 0{,}5 \times 0{,}5}{\Delta i_{L,max} \cdot f} = \frac{E}{4 \cdot \Delta i_{L,max} \cdot f}
\]

\medskip
\hrule\vspace{4pt}

% -- 1.2 Hacheur Boost (elevateur) ------------------------------------
\textbf{1.2 Hacheur Boost - élévateur}

\smallskip
En régime permanent, $\langle v_L \rangle = 0$ sur une période :
\[
  E_b \cdot \alpha + (E_b - V_{DC})(1-\alpha) = 0
  \quad\Rightarrow\quad
  \boxed{V_{DC} = \frac{E_b}{1-\alpha}} > E_b
\]
Phase $Q_1$ fermé ($t \in [0\,;\,\alpha T[$) : $v_{L} = E_b$ (courant croît).\\
Phase $Q_1$ ouvert ($t \in [\alpha T\,;\,T[$) : $v_{L} = E_b - V_{DC}$ (courant décroît, diode passe).

\medskip
\hrule\vspace{4pt}

% -- 1.3 Hacheur 4 quadrants (pont en H) ------------------------------
\textbf{1.3 Hacheur 4 quadrants - pont en H}

\smallskip
\begin{center}
\renewcommand{\arraystretch}{1.35}
\begin{tabular}{lll}
\toprule
\textbf{Configuration} & \textbf{Tension moyenne} & \textbf{Plage} \\
\midrule
Hacheur Buck (1Q) & $\langle u \rangle = \alpha \cdot E$ & $0 \leq \alpha \leq 1$ \\[4pt]
Pont en H bipolaire (4Q) & $\langle u \rangle = (2\alpha - 1) \cdot E$ & $\alpha = 0{,}5 \to u = 0$ \\[4pt]
Entrelacement $n$ bras & $\Delta i_{resultante} = \Delta i_L / n$ & Déphasage $T/n$ entre bras \\
\bottomrule
\end{tabular}
\end{center}

\smallskip
Conservation de la puissance (Boost idéal) : $E_b \cdot I_{E} = V_{DC} \cdot I_{DC}$.

\end{tcolorbox}

% -----------------------------------------------------------------------
%  2. FORMULAIRE SYNTHETIQUE
% -----------------------------------------------------------------------

\begin{tcolorbox}[
  ficheSection=colElec,
  title={2. Formulaire synthétique},
]
\small
\begin{multicols}{2}
\renewcommand{\arraystretch}{1.45}
\begin{tabular}{@{}p{4.5cm}ll@{}}
\toprule
\textbf{Formule} & \textbf{Usage} & \textbf{U.} \\
\midrule
\multicolumn{3}{@{}l}{\footnotesize\textit{Hacheur Buck}} \\
$\langle u_s \rangle = \alpha \cdot E$ & Tension sortie & V \\
$\alpha = \langle u_s \rangle / E$ & Rapport cyclique & - \\
$\Delta i_L = E\alpha(1-\alpha)/(Lf)$ & Ondulation & A \\
$L_{min} = E/(4\Delta i_{max} f)$ & Inductance min & H \\
$t_{on} = \alpha T$ & Durée passant & s \\
\midrule
\multicolumn{3}{@{}l}{\footnotesize\textit{Hacheur Boost}} \\
$V_{DC} = E_b/(1-\alpha)$ & Tension sortie & V \\
$\langle v_L \rangle = 0$ & Régime perm. & - \\
$E_b I_E = V_{DC} I_{DC}$ & Bilan puiss. & W \\
\bottomrule
\end{tabular}

\columnbreak

\renewcommand{\arraystretch}{1.45}
\begin{tabular}{@{}p{4.5cm}ll@{}}
\toprule
\textbf{Formule} & \textbf{Usage} & \textbf{U.} \\
\midrule
\multicolumn{3}{@{}l}{\footnotesize\textit{Pont en H (4Q)}} \\
$\langle u \rangle = (2\alpha-1)E$ & Tens. moy. bip. & V \\
$\alpha > 0{,}5 \to u > 0$ & Avance & - \\
$\alpha < 0{,}5 \to u < 0$ & Revers. & - \\
\midrule
\multicolumn{3}{@{}l}{\footnotesize\textit{Entrelacement}} \\
$\Delta i_{res} = \Delta i_L / n$ & Ondulation rés. & A \\
$n$ bras déphasés de $T/n$ & Condition & - \\
\midrule
\multicolumn{3}{@{}l}{\footnotesize\textit{Données typiques}} \\
$f = 50\,\mathrm{kHz}$ & Fréq. comm. & Hz \\
$\alpha_{max,\Delta i} = 0{,}5$ & Pire cas & - \\
\bottomrule
\end{tabular}
\end{multicols}
\normalsize
\end{tcolorbox}

% -----------------------------------------------------------------------
%  3. QUESTIONS TYPES
% -----------------------------------------------------------------------

\begin{tcolorbox}[
  ficheSection=colElec,
  title={3. Questions types - extraites des sujets réels},
]

% -- QT1 -------------------------------------------------------------

\begin{tcolorbox}[questiontype,
  title={QT1 - Hacheur Buck - rapport cyclique et inductance minimale \hfill\normalfont\footnotesize E2 disciplinaire 2023 (robot Bakus)}
]
\textbf{Énoncé.}
Un convertisseur dévolteur recharge une batterie Li-ion depuis une LiFePO4 ($E = 500\,\mathrm{V}$). Conduction continue, $\langle u_{sortie} \rangle \in [70\,;\,117{,}5]\,\mathrm{V}$, $\Delta i \leq 2\,\mathrm{A}$, $f = 50\,\mathrm{kHz}$.
\begin{enumerate}[label=(\alph*)]
  \item Exprimer $\langle u_{sortie} \rangle$ en fonction de $E$ et $\alpha$.
  \item Donner la plage de $\alpha$.
  \item Calculer $L_{min}$.
\end{enumerate}

\tcbline

\textbf{Corrigé.}

\textbf{Étape 1 (a) — Loi de la valeur moyenne.} Pendant $[0;\alpha T[$ l'interrupteur conduit : $u_{sortie} = E$. Pendant $[\alpha T;T[$ la diode roue-libre conduit : $u_{sortie} = 0$.
\[
  \langle u_{sortie} \rangle = \alpha \cdot E
\]

\textbf{Étape 2 (b) — Plage de $\alpha$.} $\alpha_{min} = 70/500 = 0{,}14$ ; $\alpha_{max} = 117{,}5/500 = 0{,}235$.

\textbf{Étape 3 (c) — Inductance minimale.} L'ondulation est maximale pour $\alpha = 0{,}5$ (pire cas) :
\[
  L_{min} = \frac{E}{4 \cdot \Delta i_{max} \cdot f} = \frac{500}{4 \times 2 \times 50\,000} = \frac{500}{400\,000} = \mathbf{1{,}25\,\mathrm{mH}}
\]

\textbf{Conclusion.} $L_{min} = 1{,}25\,\mathrm{mH}$ garantit $\Delta i \leq 2\,\mathrm{A}$ même au point de fonctionnement le plus défavorable ($\alpha = 0{,}5$), alors que le point réel se situe dans $[0{,}14\,;\,0{,}235]$ : la marge prise satisfait le cahier des charges sur toute la plage utile.
\end{tcolorbox}

\vspace{4pt}

% -- QT2 -------------------------------------------------------------

\begin{tcolorbox}[questiontype,
  title={QT2 - Hacheur Boost - dérivation par lois de Kirchhoff \hfill\normalfont\footnotesize Rapport jury 2025 (Q27-Q32)}
]
\textbf{Énoncé.}
Pour un hacheur élévateur ($E_b$ source, $V_{DC}$ charge) avec interrupteur $Q_1$ et diode $D_1$ : écrire les lois de Kirchhoff pendant les deux phases. En déduire $V_{DC} = E_b/(1-\alpha)$.

\tcbline

\textbf{Corrigé.}

\textbf{Étape 1 — Phase 1} ($Q_1$ fermé, $D_1$ bloquée) : $v_L = E_b > 0$ - courant $i_L$ croît.

\textbf{Étape 2 — Phase 2} ($Q_1$ ouvert, $D_1$ passante) : $v_L = E_b - V_{DC} < 0$ - courant $i_L$ décroît.

\textbf{Étape 3 — Régime permanent.} $\langle v_L \rangle = 0$ sur une période :
\[
  E_b \cdot \alpha T + (E_b - V_{DC})(1-\alpha)T = 0
  \quad\Rightarrow\quad
  E_b = V_{DC}(1-\alpha)
  \quad\Rightarrow\quad
  \boxed{V_{DC} = \frac{E_b}{1-\alpha}}
\]

\textbf{Conclusion.} La relation $V_{DC} = E_b/(1-\alpha)$ confirme que le Boost ne peut qu'élever la tension ($V_{DC} > E_b$ dès que $\alpha > 0$) : c'est ce qui justifie son choix chaque fois que le cahier des charges impose $V_{DC} > E_b$. \textit{Le rapport jury 2025 relève que de nombreux candidats ne maîtrisent pas les lois basiques de l'électricité (loi des nœuds, loi des mailles) : écrire systématiquement $v_L$, $v_D$, $v_Q$ pour chaque état.}
\end{tcolorbox}

\vspace{4pt}

% -- QT3 -------------------------------------------------------------

\begin{tcolorbox}[questiontype,
  title={QT3 - Entrelacement de 4 bras - ondulation résultante \hfill\normalfont\footnotesize E2 disciplinaire 2023 (question C7)}
]
\textbf{Énoncé.}
Le convertisseur Buck du robot Bakus est réalisé avec 4 bras entrelacés (déphasage $T/4$ entre bras), chaque bras ayant $L = 1{,}25\,\mathrm{mH}$. L'ondulation par bras est $\Delta i_L = 2\,\mathrm{A}$. Quelle est l'ondulation résultante du courant de charge ?

\tcbline

\textbf{Corrigé.}

\textbf{Étape 1 — Principe.} Avec $n = 4$ bras déphasés de $T/4$, les ondulations se compensent partiellement.

\textbf{Étape 2 — Calcul.} L'ondulation résultante est divisée par $n$ :
\[
  \Delta i_{resultante} = \frac{\Delta i_L}{n} = \frac{2}{4} = \mathbf{0{,}5\,\mathrm{A}}
\]

\textbf{Conclusion.} À inductance $L$ identique par bras, l'ondulation résultante est réduite d'un facteur 4 grâce à l'entrelacement — ou, à spécification d'ondulation égale, cela permet de réduire l'inductance par bras d'un facteur 4, ce qui allège la solution. \textit{Condition nécessaire : déphasage exactement $T/4$ entre chaque bras. Sans déphasage, les ondulations s'additionnent.}
\end{tcolorbox}

\end{tcolorbox}

% -----------------------------------------------------------------------
%  4. ATTENTES DU JURY
% -----------------------------------------------------------------------

\begin{tcolorbox}[
  jury,
  title={4. Attentes du jury - Erreurs récurrentes et points de vigilance},
]

\begin{enumerate}[label=\textbf{\arabic*.}, leftmargin=*, itemsep=5pt]

  \item \textbf{Confondre Buck et Boost.}
  Buck (dévolteur) : $V_{out} = \alpha E < E$ (rapport cyclique $< 1$).
  Boost (élévateur) : $V_{out} = E/(1-\alpha) > E$. L'erreur classique est d'appliquer la formule Buck à un schéma Boost.

  \item \textbf{Signe de $v_L$ en phase OFF du Boost.}
  $v_L = E_b - V_{DC} < 0$ : le courant \textbf{décroît}. Beaucoup de candidats obtiennent un courant croissant dans les deux phases en inversant le signe.

  \item \textbf{La condition $\langle v_L \rangle = 0$ est la clé.}
  Sans cette condition de régime permanent, il est impossible de dériver les formules de tension moyenne. L'énoncer explicitement avant tout calcul.

  \item \textbf{Ondulation maximale à $\alpha = 0{,}5$, pas à $\alpha$ nominal.}
  Calculer $L_{min}$ pour $\alpha = 0{,}5$ (pire cas), même si le point de fonctionnement nominal est $\alpha = 0{,}14$.

  \item \textbf{Identifier la nature du convertisseur depuis le schéma.}
  Le jury apprécie de nommer Buck/Boost/4Q et de justifier par le sens du flux d'énergie (sens courant, rapport $V_{out}/V_{in}$).

  \item \textbf{Entrelacement : déphasage obligatoire.}
  L'ondulation est divisée par $n$ \textbf{seulement si} les bras sont déphasés de $T/n$.

\end{enumerate}

\medskip
\begin{tcolorbox}[
  colback=colJuryFrame!12!colJury,
  colframe=colJuryFrame,
  left=4pt, right=4pt, top=2pt, bottom=2pt,
  fontupper=\small\itshape
]
Rapport jury 2025 (commentaires du jury, épreuve disciplinaire) : «~De nombreux candidats ne maitrisent pas les lois basiques de l'électricité (la loi des n\oe uds, la loi des mailles, les calculs de puissance ou encore la loi d'Ohm).~»\\
Le corrigé du rapport jury 2025 (Q14) identifie explicitement CV1 = hacheur, CV2 = redresseur, CV3 = onduleur, chacun avec sa fonction de conversion : nommer chaque convertisseur et justifier sa fonction est un attendu de base.
\end{tcolorbox}

\end{tcolorbox}

% -----------------------------------------------------------------------
%  PIED - PRESENCE DANS LES SUJETS
% -----------------------------------------------------------------------

\begin{tcolorbox}[
  ficheSection=colPresence,
  title={Présence dans les sujets analysés (2013-2026)},
]
\begin{center}
\footnotesize\renewcommand{\arraystretch}{1.15}
\begin{tabular}{lccccccccccccc}
\toprule
\textbf{Année} & 13 & 14 & 15 & 16 & 17 & 18 & 19 & 20 & 21 & 22 & 23 & 24 & 25 \\
\midrule
E1 (commune)   & - & - & $\checkmark$ & - & $\checkmark$ & - & - & $\checkmark$ & - & $\checkmark$ & - & - & $\checkmark$ \\
E2 discipl. IE & $\checkmark$ & - & - & $\checkmark$ & - & $\checkmark$ & $\checkmark$ & - & $\checkmark$ & $\checkmark$ & $\checkmark$ & $\checkmark$ & $\checkmark$ \\
\bottomrule
\end{tabular}
\end{center}
\smallskip
\footnotesize\textbf{Sous-thèmes fréquents :}
hacheur Buck - rapport cyclique (8/13) - inductance minimale (5/13) - Boost / MPPT (4/13) - hacheur 4 quadrants / pont en H (3/13) - entrelacement de bras (2/13)

\medskip
\textbf{Ordres de grandeur (repères rapides).}
\begin{itemize}[leftmargin=1.3em]
  \item Réseau BT FR : tension phase-neutre typique \SI{230}{\volt}, tension phase-phase typique \SI{400}{\volt}, fréquence \SI{50}{\hertz}.
  \item Courants d'étude fréquents en sujet : de quelques \si{\ampere} à quelques dizaines de \si{\ampere}.
  \item En conversion d'énergie, un rendement attendu reste généralement entre 0{,}8 et 0{,}98 selon la technologie.
\end{itemize}

\medskip
\textbf{Synthèse express (à refaire sans support).}
\begin{enumerate}[leftmargin=1.4em]
  \item Poser les hypothèses du modèle et les conventions de signe.
  \item Écrire 2 à 4 relations clés avec unités.
  \item Réaliser un calcul numérique complet et vérifier l'ordre de grandeur.
  \item Justifier le résultat par un argument technique de niveau CAPET.
\end{enumerate}

\end{tcolorbox}

\CAPETENDFILE
